\documentclass[11pt]{amsart}
\usepackage[expansion=false]{microtype}
\usepackage[a4paper,textwidth=149mm,textheight=235mm,centering]{geometry}
\usepackage{amsmath,amssymb,mathtools,booktabs,enumitem}
\usepackage[colorlinks,linkcolor=blue!50!black,citecolor=blue!50!black,urlcolor=blue!50!black]{hyperref}
\usepackage{xcolor}
\setlength{\emergencystretch}{1.4em}
\setcounter{tocdepth}{1}
\numberwithin{equation}{section}
\newtheorem{theorem}{Theorem}[section]
\newtheorem{lemma}[theorem]{Lemma}
\newtheorem{proposition}[theorem]{Proposition}
\newtheorem{corollary}[theorem]{Corollary}
\newtheorem{claim}[theorem]{Claim}
\newtheorem*{mainthm}{Theorem}
\DeclareMathOperator{\lcm}{lcm}
\DeclareMathOperator{\Tr}{Tr}
\DeclareMathOperator{\FPdim}{FPdim}
\DeclareMathOperator{\Irr}{Irr}
\DeclareMathOperator{\Gal}{Gal}
\newcommand{\QQ}{\mathbb Q}
\newcommand{\ZZ}{\mathbb Z}
\newcommand{\C}{\mathcal C}
\newcommand{\one}{\mathbf1}
\newcommand{\supfile}[2]{\href{https://github.com/sebastienpalcoux/WGT-fusion-categories/blob/433d1c45795202991f1693e00d9770e4cccf05f6/papers/rank-bound-supplement/#1}{#2}}
\allowdisplaybreaks[2]
\title[A rank bound for simple integral modular categories]{A rank bound for simple integral modular categories}
\author{Sebastian Burciu}
\author{Jingcheng Dong}
\author{S\'ebastien Palcoux}
\date{}
\hypersetup{pdftitle={A rank bound for simple integral modular categories},pdfauthor={Sebastian Burciu, Jingcheng Dong and Sebastien Palcoux}}
\begin{document}
\begin{abstract}
Every non-pointed simple integral modular category has rank at least $32$.
The sum of the squared simple dimensions becomes a reciprocal-square equation
whose denominators can be enumerated from the rank alone. Character
orthogonality and the local modular-group action restrict the multiplicities
of small denominators. Exact tensor-product and character tests eliminate the
remaining possibilities through rank $31$.
\end{abstract}
\maketitle
\section{The dimension equation}\label{sec:dimensions}
A fusion category over $\mathbb C$ is \emph{simple} if its only fusion
subcategories are itself and $\mathrm{Vec}$, \emph{pointed} if every
simple object is invertible, and \emph{integral} if every simple object
has integer Frobenius--Perron dimension. The rank counts the isomorphism
classes of simples. A modular category carries a spherical structure and
a nondegenerate braiding. Throughout, the ground field is $\mathbb C$.
\begin{mainthm}
Every non-pointed simple integral modular category has rank at least $32$.
\end{mainthm}

Let $d_i$ be the simple dimensions and $D=\sum_i d_i^2$. Although $D$ is
unknown, each $d_i^2$ divides $D$. On writing $D=qM^2$ with $q$ square-free,
the integers $M/d_i$ satisfy a reciprocal-square equation. The rank bounds
its least denominator, and each initial segment bounds the next denominator.
The main difficulty is a long run of equal small denominators. It represents
a whole class of simples of the same dimension, whose characters and twists
constrain the other classes. For example, in square dimension we exclude
denominator three and bound the multiplicity of four by thirteen through
rank $31$. The local restrictions in Section~\ref{blk:section} give the starting
codegrees in Section~\ref{sec:small} and the bounds on large classes in
Section~\ref{sec:large}. They apply without fixing a family of prime
factorizations of $D$, so the rank theorem can be read independently
of the dimension results in \cite{BDP}. The only computer-assisted step
is the exhaustive finite check in Section~\ref{sec:finish}.
The \supfile{README.md}{supplement guide} explains how to reproduce it.

\subsection{Character identities}
Use the positive spherical structure: integrality implies pseudo-unitarity
\cite[Proposition~8.24]{ENOannals}, and nondegeneracy depends only on the
braiding. The pointed subcategory is proper, so the unit $0$ is the only
invertible simple. Normalize $S$ by $S_{0i}=d_i$. Then $S=S^t$,
$S\overline S^t=DI$, and $\chi_i(j)=S_{ij}/d_i$ are the fusion characters.
Their values are cyclotomic integers, being eigenvalues of integral fusion
matrices in the modular field. The \emph{formal codegree}
$f_i=\sum_j|\chi_i(j)|^2=D/d_i^2$ is a rational algebraic integer, hence an
integer. Galois permutes the characters preserving $f_i$ and $d_i$. By
symmetry, $\sigma\chi_i(j)=\chi_{\sigma i}(j)=\chi_i(\sigma j)$.
On $s=S/\sqrt D$ this action has the common sign
$\sqrt D/\sigma(\sqrt D)$. Write $T=\operatorname{diag}(\theta_i)$ for
the twist matrix, with $\theta_{i^*}=\theta_i$. The Gauss sum
$p^+=\sum_i d_i^2\theta_i$ has modulus $\sqrt D$. Its phase
$\xi=p^+/\sqrt D$ is the multiplicative central charge. In this
normalization the modular relation is $sTs=\xi T^{-1}sT^{-1}$.

\begin{lemma}\label{lem:basic-table}
Every nonunit row of $S$ has a zero, and $S_{ij}=0$ for coprime nonunit
degrees $d_i,d_j$. Each nonunit degree has at least two distinct prime
factors. Every prime dividing $D$ occurs to exponent at least two.
\end{lemma}
\begin{proof}
Put $s_0=d_i$, $t=d_j$. Both $S_{ij}/s_0$ and $S_{ij}/t$ are integral,
so B\'ezout gives integrality of $z_{ij}=S_{ij}/\lcm(s_0,t)$.
A nonzero algebraic integer has Galois mean square at least one by the
norm and the arithmetic--geometric mean inequality. A nonvanishing row
$i\ne0$ would therefore give $D\ge\sum_j\lcm(s_0,d_j)^2>D$ upon averaging
its norm: conjugation commutes with Galois, and the unit gives strictness.

Each conjugate $\chi$ of $\chi_i$ has norm square $D/s_0^2<D$, and
$|\chi(j)|\le d_j$ by the spectral-radius bound. Equality passes to
duals and, by the triangle inequality, to every constituent of a product
of labels attaining equality. A nonunit label generates the category, so
equality there would force norm square $D$. Thus all conjugates of
$S_{ij}/(s_0t)$ have modulus less than one for $i,j\ne0$. If
$\gcd(s_0,t)=1$, this is $z_{ij}$, whose integral norm forces it to vanish.
If $s_0$ were a prime power $p^a$, every term in
$\sum_{j\ne0}d_j\chi_i(j)=-1$ would be divisible by $p$.

Finally $P=\prod_j[X_j][X_j^*]$ vanishes at every character except the
dimension character. Diagonalizing by $S$ gives its unit coefficient
$\prod_j d_j^2/D\in\mathbb Z$. A prime dividing $D$ only once could divide
no $d_j$, since $d_j^2\mid D$, contradicting this divisibility.
\end{proof}

At least three primes divide $M$. Indeed $d_i\mid M$ for every $i$.
With at most two primes, all nonunit degrees would share a prime,
contradicting orthogonality with the unit.

\subsection{Enumerating the denominators}
\begin{lemma}\label{lem:deletion}
For each prime $p$, at least two nonunit labels have degree prime to $p$.
\end{lemma}
\begin{proof}
Reducing $\sum_{j\ne0}d_j\chi_i(j)=-1$ above $p$ gives a $p$-free nonunit
column meeting row $i$, for each $i\ne0$. A unique such column would have
no zero, contrary to Lemma~\ref{lem:basic-table} and symmetry.
\end{proof}

Order the nonunit denominators $b_i=M/d_i$ as $b_1\le\cdots\le b_{r-1}$.
Since $M>b_1$,
\begin{equation}\label{eq:recip}
 q=M^{-2}+\sum_{i=1}^{r-1}b_i^{-2},\qquad q b_1^2<r .
\end{equation}
After a nonempty prefix put $L=\lcm(b_1,\dots,b_j)$, $h=r-1-j$ and
$E=L^2(q-\sum_{i\le j}b_i^{-2})$. If $h>0$, the next denominator $b$ obeys
\begin{equation}\label{eq:interval}
 b\ge b_j,\qquad L^2<Eb^2\le(h+1/36)L^2 .
\end{equation}
Indeed the remainder exceeds $b^{-2}$ and is at most $h/b^2+1/(6b)^2$,
since $M/b\ge6$. Setting $g=\gcd(L,b)$ and $k=b/g$ updates $(L,E)$ to
$(kL,k^2E-(L/g)^2)$. This gives a finite tree from the rank alone.

With two positions left, write their degrees $d\le e$ and $M=kL$. All
earlier nonunit degrees are multiples of $k$. Lemma~\ref{lem:deletion}
forces $\gcd(k,de)=1$, hence $d,e\mid L$. The exact closure is
$Ek^2=1+d^2+e^2$ with $M/e\ge b_j$. One can pair divisors of $L$ by
squared residues modulo $E$, or enumerate $b=M/e$ in \eqref{eq:interval}:
then $g=\gcd(L,b)=L/e$, $k=b/g$ and $d^2=Ek^2-1-e^2$. Both give every
completion satisfying the degree restrictions. Reduction modulo four and
five gives $\gcd(k,10)=1$, whence $E\not\equiv0,4,7\pmod8$.
With one position left, the same lemma gives $k=1$ and $E=1+d^2$.

For an example of what this arithmetic allows, take $M=140$, $q=1$.
The denominators $4,5,7,10$ with respective multiplicities $11,4,6,3$
give the nonunit degrees $35,28,20,14$. Together with the vacuum they
satisfy \eqref{eq:recip}. We shall see below how one local character
condition eliminates this solution.

\subsection{The field of a row}
\begin{lemma}\label{conduct}
The field $K_i=\QQ(\chi_i(j):j)$ has conductor dividing $f_i$.
\end{lemma}
\begin{proof}
Write $f=f_i$. For $a$ in the inverse different of $K_i$, the fusion-ring
element $\sum_j\Tr(a\chi_i(j^*))[j]$ has eigenvalue $fa$ at $i$ by
orthogonality. Hence the different divides $f\mathcal O_{K_i}$.

At $p$ absorb the unramified part into the base $F$, so the remaining
extension $E/F$ lies in $L=F(\zeta_{p^v})$ of least possible level.
Its normalized different exponent is at most $v_p(f)$. For odd $p$ write
$[E:F]=e_0p^s$, $e_0\mid p-1$. In the ramified case cyclicity of the
wild part gives $v=s+1$ and $L/E$ is tame. From $\Phi'_{p^v}$ and
transitivity, the normalized exponent is
$s+1-(1-p^{-s})/(p-1)-1/(e_0p^s)$. It is positive for tame
ramification and exceeds $s$ otherwise, so $p^v\mid f$.
At two, conductor four has normalized exponent one. For $v\ge3$, the
stabilizer of $E$ misses $1+2^{v-1}$, meets $\langle5\rangle$ trivially,
and has order at most two. If nontrivial it sends $\zeta$ to
$\pm\zeta^{-1}$, whose difference from $\zeta$ has valuation two.
The normalized different of $E/F$ is at least
$v-1-2/2^{v-1}>v-2$. Thus the conductor divides $2f$.

Put $m=2^{v_2(f)}$. If $m=1$, the extra factor two changes no field.
Otherwise take $F$ to be an odd cyclotomic field containing the odd part,
and $\zeta$ of order $2m$. Suppose $e=1+m$ moves row $i$.
Write $\chi_i(Y)=\sum_{k=0}^{m-1}z_{kY}\zeta^k$ with
$z_{kY}\in\mathcal O_F$ (the conductors are coprime). For
$G=(\ZZ/2m)^\times$ and the row stabilizer $H$, Fourier inversion gives
$z_{1Y}=m^{-1}\sum_{g\in G}\zeta^{-g}\chi_{gi}(Y)$. Orthogonality yields
\[
 \sum_Y|z_{1Y}|^2=\frac f{m^2}\sum_{h\in H}\sum_{g\in G}\zeta^{g(h-1)}
 =\frac f m ,
\]
as only $h=1,e$ can contribute and $e\notin H$. Yet $e$ permutes columns
and negates $z_{1Y}$, so the nonzero coefficients pair off. The left side
is divisible by two in the algebraic integers, whereas $f/m$ is odd.
\end{proof}

\subsection{The principal block}\label{blk:principal}
Let $R=K_0(\C)$ and $e_i$ be the primitive idempotent of $\chi_i$.
The unit coefficient $\tau(a)=\sum_i\chi_i(a)/f_i$ makes $R$ its own
trace dual, also locally, since $\tau([i][j]^*)=\delta_{ij}$.
At a place above $p\mid D$, the \emph{principal block} $P$ consists of the
rows satisfying $\chi_i(j)\equiv d_j$ for every $j$. Galois permutes
argument labels within degree classes, so this holds at every place above
$p$ and $P$ is Galois stable. Moreover $e_P=\sum_{i\in P}e_i\in
R\otimes\ZZ_p$: reduce this algebra modulo $p$ and lift the idempotent
of the local factor for the dimension character. Its residue field is
$\mathbb F_p$, and the lifted idempotent has value one precisely on $P$.

\begin{lemma}\label{blk:support}
For every $i$, $\sum_{j\in P}|\chi_i(j)|^2$ is an integer divisible by
$p^{v_p(f_i)}$. If $f_i$ is a $p$-power, its row is supported in $P$,
with equality of supports if $p\nmid d_i$.
\end{lemma}
\begin{proof}
The sum is integral and Galois invariant. Divided by $f_i$ it is
$\tau(e_P[i][i]^*)\in\ZZ_p$ by symmetry. It lies between the vacuum
contribution one and $f_i$, proving the support bound. For $j\in P$,
$\chi_j(i)\equiv d_i$ gives a nonzero entry when $p\nmid d_i$.
\end{proof}

\begin{lemma}\label{blk:height}
If $p^2\mid D$, then $p^2\mid f_i$ for each $i\in P$.
\end{lemma}
\begin{proof}
The unit is clear. For another $i\in P$ let $O$ be its Galois orbit,
$h=|O|$, $e_O=\sum_{j\in O}e_j$. For $a\in R$, the integer
$t(a)=\sum_{j\in O}\chi_j(a)$ is congruent to $h\chi_0(a)$ modulo $p$.
Thus $b=(hD e_0-f_i e_O)/p$ belongs to $R$ by trace duality, since
$\tau(ba)=(h\chi_0(a)-t(a))/p$. Now
$-f_i/p=\chi_i(b)\equiv\chi_0(b)=hD/p\equiv0\pmod p$.
\end{proof}

\begin{lemma}\label{blk:witness}
Let $A_p$ consist of the degrees $d_j>1$ such that $p^2\mid f_j$ and
$\gcd(d_j,d_i)>1$ for every $p$-free nonunit degree $d_i$. If $p^2\mid D$,
then $\gcd(A_p)=1$.
\end{lemma}
\begin{proof}
Each nonunit $j\in P$ contributes to $A_p$ by Lemma~\ref{blk:height}:
$\chi_j(i)\equiv d_i$ gives a nonzero crossing for every $p$-free $i$.
Categorical Cauchy gives a label $x$ with nontrivial $p$-part of its twist.
Block orthogonality (Lemma~\ref{blk:local}) gives
$0=\sum_{j\in P}d_j\chi_x(j)=1+\sum_{0\ne j\in P}d_j\chi_x(j)$.
No prime divides every degree in the latter sum.
\end{proof}

Return to the example $M=140$. The only nonunit degrees with
$4\mid D/d^2$ are $14$ and $35$. Both are divisible by seven, so
Lemma~\ref{blk:witness} at two already rules out the list.

\section{What a prime in the dimension forces}\label{blk:section}
Let $N=\operatorname{ord}(T)$. The congruence theorem gives a unitary
representation of $\mathrm{SL}_2(\ZZ/L)$ sending the Weyl matrix
$W=\left(\begin{smallmatrix}0&-1\\1&0\end{smallmatrix}\right)$ to $s$ and
the upper unipotent $U=\left(\begin{smallmatrix}1&1\\0&1\end{smallmatrix}\right)$
to $T/\gamma$, with $\gamma$ of two-power order and $L=\operatorname{ord}(T/\gamma)$
\cite[Theorems II and 6.10]{DLNCongruence}. Chinese remaindering gives
commuting local factors. Write $w=W_p$, $u=U_p$, and
$h_a=\operatorname{diag}(a,a^{-1})_p$. The entries of $u$ are the
$p$-parts $\theta_{i,p}$ of the twists up to a common scalar, which is one
for odd $p$. The torus $h_a$ acts by a signed Galois permutation preserving
degrees and raising the normalized local twists to their $a^2$th powers
(up to replacing $a$ by $a^{-1}$). Its sign at a fixed label is that of
the Galois action on $d_i/\sqrt D$. By categorical Cauchy, $N$ and $D$
have the same prime divisors \cite[Theorem~3.9]{BNRWFinite}. Here
$\gamma^3=\xi$. We call $\theta_{i,p}$ the $p$-phase of $i$, and call
$i$ active at $p$ if this phase differs from one. A phase fiber consists
of all labels with the specified primary phases. Its squared mass is the
sum of their squared degrees.

\begin{lemma}\label{blk:local}
Each full $p$-block spans an invariant space for the $p$-local group.
Equivalently, if $B$ is a block and $\theta_{x,p}\ne\theta_{y,p}$, then
$\sum_{z\in B}s_{xz}\overline{s_{yz}}=0$.
\end{lemma}
\begin{proof}
Blocks are the fibers of reduction of the characters at a place
$\mathfrak p$ above $p$. Let $L_j^i$ be the $j$th generalized rotation
trace with central insertion $(i,i^*)$ in $\C\boxtimes\C^{\rm rev}$,
which has twist one. Put $A_j=sT^js^{-1}$. The indicator formula and
strand grouping \cite[Lemmas~2.5, 2.7, Proposition~2.8, Corollary~5.5
and Proposition~6.1]{NSind} give
\[
 L_j^i(e_z)=|(A_j)_{iz}|^2,\qquad
 L_j^i(y)^p\equiv L_{j/p}^i(y^p)\pmod{\mathfrak p}\quad(p\mid j).
\]
Indeed, on $[c]$ the indicator sums
$s_{ia}s_{ib}(\theta_a/\theta_b)^jN_{ab}^c$ over $a,b$. Contracting with
$f_z e_z=\sum_c\overline{\chi_z(c)}[c]$ gives conjugate factors. For the
congruence, the $p$th power of a trace reduces to the trace of the $p$th
power of the rotation, which groups $p$ strands. Extend $L_j^i$ linearly;
additivity and Frobenius give the congruence for integral cyclotomic
coefficients.

Let $m=N_{p'}$ be the prime-to-$p$ part of $N$. Choose $b\ge\max(1,v_p(N))$
with $p^b\equiv1\pmod m$. For an object $A$ and $n\ge2b$,
\[
 L_m^i(A^{p^n})-L_{mp^b}^i(A^{p^{n-b}})\in
 p^{n-2b+1}\overline{\ZZ}_p .
\]
These are traces of the powers $p^n,p^{n-b}$ of one rotation on $mp^n$
strands. Its eigenvalues have orders dividing $mp^n$, with multiplicities
constant on Galois orbits over $\QQ(\zeta_N)$ since all power traces lie
there \cite[Proposition~5.7]{NSind}, \cite[Theorem~7.7]{NSfs}. The two
powers agree on prime-to-$p$ parts and on roots of order dividing
$p^{n-b}$. For exact $p$-part $p^j>p^{n-b}$, units congruent to one
modulo $p^b$ and trivial away from $p$ give orbits of size $p^{j-b}$.
Both powers are constant on each orbit, giving the divisibility.

A maximal ideal of the fusion ring modulo $p$ with residue field
$\mathbb F_{p^h}$ gives $h$ Frobenius-conjugate blocks $B_j$. For each
nonzero residue $a$, lift the element with value $a$ there and zero at
the other maximal ideals to an object $A$, taking nonnegative basis
coefficients modulo $p$. Choose $h\mid b,n$. In character coordinates
both powers of $A$ then tend as $n\to\infty$ to
$\sum_j[a]^{p^j}e_{B_j}$, where $[a]$ is the Teichm\"uller lift and
$e_B=\sum_{z\in B}e_z$. The two trace forms agree on these sums, hence
on each $e_B$ by independence of the characters $a\mapsto[a]^{p^j}$.
Since $N\mid mp^b$, the matrix formula gives
$\sum_{z\in B}|(A_m)_{iz}|^2=\mathbf1_{i\in B}$. Thus $A_m$ preserves
$B$ by unitarity. As $T^m$ distinguishes the $p$-phases, the block
projector commutes with $sus^{-1}$ and with $u$, the two generating
unipotents. Conjugating the projector gives the stated orthogonality
($s^2$ is charge conjugation and commutes with twists).
\end{proof}

\begin{lemma}\label{blk:admissible}
If $w$ preserves the line of $z$, then $p\nmid f_z$. Row $z$ meets a
column of nontrivial $p$-twist exactly when $p\mid f_z$. A nontrivial
$p$-twist at $z$ requires $p\mid f_z$ and, for odd $p$, its order divides $f_z$.
\end{lemma}
\begin{proof}
Assume $p\mid D$. Choose $k$ divisible by $p$ and the $p$-part of $L$,
with $k\equiv1\pmod{L_{p'}}$. On the complementary factors,
$U(WUW^{-1})U=W^{-1}$ identifies $sT^ks^{-1}$ up to diagonal phases
with the inverse complementary Weyl. If $w$ preserves the line of $z$,
its inverse acts there by a phase. The complementary inverse Weyl
therefore has $(0,z)$ entry of modulus $|(s^{-1})_{0z}|=|s_{0z}|$.
Hence $L_k^0(f_ze_z)=1$. The preceding congruence at $y=f_ze_z$ gives
$1\equiv f_z^{p-1}L_{k/p}^0(y)\pmod{\mathfrak p}$, so $p\nmid f_z$.

If $p\nmid f_z$, the character expansion makes $e_z$ locally integral:
$z$ forms a singleton block and its row misses all active columns by
Lemma~\ref{blk:local} with $x=0$. Conversely, if it misses them,
$s\operatorname{diag}(\theta_{i,p})s^{-1}$ fixes $v_z$. Both generating
unipotents preserve its line, as does $w$. There, $sus^{-1}=wuw^{-1}$
must act by the same scalar as $u$, giving $\theta_{z,p}=1$.

If the odd local twist has order $p^j$, every automorphism fixing the row
has multiplier $\pm1\pmod{p^j}$ by square Galois covariance. The row
field contains $\QQ(\zeta_{p^j}+\zeta_{p^j}^{-1})$, of conductor $p^j$
except at $p^j=3$. Now use Lemma~\ref{conduct}, and the preceding
paragraph for the exception.
\end{proof}

\begin{lemma}\label{blk:moments}
Fix an odd $p\mid D$. In a full $p$-block or in the whole label set, fix
any collection of complementary primary twists. Weight each remaining
label by $d_i^2$. If $v_p(D)$ is even, the weights at the primitive
$p$th roots are equal. If it is odd, the total weight at order $p$ equals
$p-1$ times the weight at the trivial phase.
\end{lemma}
\begin{proof}
The local factor commutes with the projector $E$ of this set by
Lemma~\ref{blk:local} and commutation of local factors. It also commutes
with $Q=sEs^{-1}$, since $s$ normalizes it. For a unit $a$ modulo the
local level, the Bruhat identity gives
\[
 D^{-1}\sum_{i\in E}d_i^2\theta_{i,p}^a
 =(Qwu^aw^{-1})_{00}=(Qu^{-1/a}wh_{-a}u^{-1/a})_{00}
 =\varepsilon(-a)(Qw)_{00}.
\]
Indeed $u$ fixes the vacuum and $h_{-a}$ acts there by
$\varepsilon(-a)$. On $\sqrt D$, pure $p$-inertia acts trivially at even
valuation and by the quadratic character at odd valuation, by the
quadratic Gauss sum. In the even case the moment is rational. Its
normalized trace to $\QQ(\zeta_p)$ removes terms of order greater than
$p$, and the only rational relation among $p$th roots is their sum.
In the odd case the average over units vanishes; roots of order $1,p,p^j$
($j\ge2$) have respective means $1,-1/(p-1),0$.
\end{proof}

Put $h_p=(p-1)/2$ for odd $p$. At each degree, phases of exact order
$p^j$ have constant multiplicity on square-unit orbits of size
$p^{j-1}h_p$, even with complementary twists fixed. We count in
\emph{packets} of $h_p$ labels of one degree.

\begin{proposition}\label{blk:packets}
For the odd primes $A$ dividing $q$, put $P=\prod_{t\in A}(h_t+1)$ and
\[
 B(q)=\max\bigl(\{2^{\mathbf1_{2\mid q}}P\}\cup
 \{(h_t+2)P/(h_t+1):t\in A\}\bigr).
\]
Every nonempty fiber fixing odd-primary components outside $A$ has at
least $B(q)$ labels. For any set $J$ of odd primes dividing $D$ to even
valuation, at least $B(q)\sum_{p\in J}(p-1)$ labels are active at some $p\in J$.
\end{proposition}
\begin{proof}
In any nonempty complementary fiber $F$, each $t\in A$ gives at least
$h_t+1$ phases and $h_t+2$ labels. Indeed a higher-order phase gives at
least $th_t\ge h_t+2$ phases. Otherwise Lemma~\ref{blk:moments} forces one
inactive label and $h_t$ primitive phases. Equality at $h_t+1$ labels
would give one inactive degree $b$ and one degree $a$ on the primitive orbit,
with $h_ta^2=(t-1)b^2$, impossible. We also have
\begin{equation}\label{blk:fibergcd}
 t\mid\sum_{i\in F}(d_i/g)^2,\qquad g=\gcd_{i\in F}d_i :
\end{equation}
the inactive and exact-order-$t$ masses sum to a multiple of $t$ after
division by $g^2$, and higher-order multiplicities are multiples of $t$.

Refining at the primes of $A$ gives at least $h_t+1$ children at each
$t$, and $h_t+2$ labels at the last prime. If $2\mid q$, all codegrees
are even. Every final odd-primary fiber is dyadically invariant, so it has
at least two labels by Lemma~\ref{blk:admissible}. This gives $2P$.

Project to $J$ and group phase tuples by their exact orders. In a
nonempty group write $J'$ for the active coordinates. With the other
phases fixed, each $p\in J'$ supplies at least $p-1$ phases of its
order, by Lemma~\ref{blk:moments} at order $p$ and the square action
above it. For $J'\ne\varnothing$ there are at least
$\prod_{p\in J'}(p-1)\ge\sum_{p\in J'}(p-1)$ tuples, each with at least
$B(q)$ labels. Every $p\in J$ occurs by Cauchy.
\end{proof}

A degree $d$ is eligible at $p$ only if $p\mid D/d^2$. At even valuation,
if only two packets fit in the eligible degree classes, higher orders are
impossible. Each primitive square class uses one packet, and equality
of squared masses forces the same degree for both.

\begin{lemma}\label{blk:activity}
For the least nonunit degree $x$ and $p\mid D$, put
$V_p=\sum_{\theta_{i,p}\ne1}d_i^2$. Then $2px V_p\ge(p-1)D$.
\end{lemma}
\begin{proof}
Put $k=N_{p'}$, $a_i=\theta_i^k$. For a simple $X$ of degree $x$, the
indicator \cite[Theorem~7.5]{NSfs} is
$\nu_k(X)=D^{-1}\sum_{i,j}N_{ij}^X d_i d_j a_i/a_j\in\ZZ$: it is
integral and lies both in the $p$-power field of the phases and in the
prime-to-$p$ field of the rotation. The weights sum to $Dx$. Equality
$\nu_k(X)=x$ would force agreement of phases along each edge of
multiplication by $X$, by reciprocity $N_{ij}^X=N_{i^*X}^{j}$ and
invariance of twists under duality. The underlying undirected graph is
connected since $X,X^*$ generate, so every phase would be trivial,
contrary to Cauchy. Thus $x-\nu_k(X)\ge1$. Averaging over Galois,
an unequal pair contributes at most $p/(p-1)$ times its normalized
weight to this deficit. It has an active index, and each choice of
active index has total weight $xV_p$ by reciprocity.
\end{proof}

At odd $p$ of even valuation let $h=(p-1)/2$. In each eligible degree
$d>1$ of multiplicity $m_d$, put $n_d=\lfloor m_d/h\rfloor$ and let $k_d$
be the difference in exact-order-$p$ packets between the two square
classes. Then $|k_d|\le n_d$ and $\sum_d k_dd^2=0$ by
Lemma~\ref{blk:moments}. After $|k_d|$ packets, opposite pairs use two
slots each; a higher order uses a multiple of $p$ slots and requires
$p^2\mid D/d^2$. Thus $V_p\le h\sum_d(n_d-\epsilon_d)d^2$, maximized
over the stated integer constraints. Here $\epsilon_d=1$ if $n_d-|k_d|$
is odd and either less than $p$ or $p^2\nmid D/d^2$; otherwise it is zero.

\begin{lemma}\label{blk:two}
If a local factor preserves the coordinate plane of $X,Y$ but its Weyl
moves the line of $X$, then $p=2$, $d_X=d_Y$, and the primary twist ratio
has order four.
\end{lemma}
\begin{proof}
The eigenvalues of $u$ on the plane are distinct, since otherwise both
generating unipotents would be scalar there. Any invariant line would
be a coordinate line with invariant orthogonal complement, so the plane
is irreducible and $w^2$ is scalar. Up to phase $w$ is a Hermitian
involution $\left(\begin{smallmatrix}a&b\\\bar b&-a\end{smallmatrix}\right)$
with $a\in\mathbb R$, $0<|b|^2=1-a^2$. Since $(wu)^3=w^2$ and $wu$ is
nondiagonal, it has projective order three and trace of modulus one.
Writing $\omega$ for the eigenvalue ratio of $u$ gives $a^2|1-\omega|^2=1$.

The complementary Weyl commutes with $u$. Both columns of $s=w_{p'}w$
meet the vacuum, so only the eigenline with its $u$-eigenvalue contributes
to their vacuum entries. Thus $d_X/d_Y=|a/b|$ or its inverse. Hence $a^2$
is rational, as is the algebraic integer $\omega+\omega^{-1}$. The possible
nontrivial prime-power orders are $2,3,4$, giving $a^2=1/4,1/3,1/2$.
Only the last yields a rational degree ratio.
\end{proof}

If $2\mid q$ and odd $p\mid M$ has even valuation, put $h=(p-1)/2$.
Then $\sum_d(\lfloor m_d/h\rfloor+\lfloor m_d/(2h)\rfloor)\ge6$ over
eligible nonunit degrees. Every nontrivial phase fiber has at least two
labels by Lemma~\ref{blk:admissible}, and a two-label fiber has equal
degrees by Lemma~\ref{blk:two}. Higher order gives at least $2ph$ labels.
Otherwise each primitive square class supplies three packets, or two
of one degree, also counted in the second floor.

\section{Starting the search}\label{sec:small}
For a nonunit $X$ of degree $n$ and codegree $f$, put $u_Y=\chi_X(Y)$,
let $c_Y$ be the Galois mean of $|u_Y|^2$, and let $H$ be its nonunit
support. The vacuum column of $STS=p^+T^{-1}ST^{-1}$ gives
$|\sum_Yd_Y\theta_Yu_Y|=\sqrt D$. Averaging and applying Cauchy--Schwarz,
then using lcm integrality and the norm on each nonzero entry, gives
\begin{equation}\label{st}
 \sqrt f\le n^{-1}\sum_Yd_Y\sqrt{c_Y},\qquad
 \sum_Yc_Y=f,\qquad c_Y\ge(d_Y/\gcd(d_Y,n))^2\quad(Y\in H).
\end{equation}
Orthogonality with the unit gives $\gcd_{Y\in H}d_Y=1$.

\begin{lemma}\label{small}
Every nonunit codegree exceeds four, and none is prime.
\end{lemma}
\begin{proof}
If $f\le4$, all supported degrees divide $n$ and $2\le|H|\le3$.
For two, both degrees are proper divisors by their gcd. Thus
$\sqrt f\le1/n+\sqrt{2(f-1)}/2$, impossible for $f=3,4$ and $n\ge6$.
For three, $f=4$ and all $c_Y=1$ on $H$. If all degrees are proper,
$2\le1/n+3/2$. At most one equals $n$, and then the other two cannot both
equal $n/2$. Thus $2n\le1+n+n/2+n/3$, forcing $n\le6$; six has no
eligible proper divisor.

If $f=p$ is prime, $D=pn^2$ is square-full, so $p\mid n$ and all
degrees divide $n$. The row is supported in the principal $p$-block
where $p^2\mid f_Y$, hence its degrees divide $n/p$. As $c_Y\ge1$
on the support, \eqref{st} gives $\sqrt p\le p^{-1}\sum_Yc_Y=1$.
\end{proof}

\begin{lemma}\label{six}
Six is not a formal codegree.
\end{lemma}
\begin{proof}
Suppose $D=6n^2$. Every degree divides $n$, $6\mid n$ by square-fullness,
and $n\ge30$ since three primes divide $n$. The entries lie in
$\ZZ[\zeta_3]$ by Lemma~\ref{conduct}. Their nonzero nonunit squared
moduli are Eisenstein norms summing to five: $(4,1)$, $(3,1,1)$, or five
ones. For $(4,1)$ the degrees are coprime divisors of $n$, both at least
six, hence both at most $n/6$, contrary to \eqref{st}. For $(3,1,1)$
all three ratios $d_Y/n$ are at most $1/2$, or one is one and the other
two are coprime divisors of $n$ at most $n/6$. In either case the
right side of \eqref{st} is at most $1/n+\sqrt{5(1+2/36)}<\sqrt6$.

For five unit entries set $k_Y=n/d_Y$ and let $m$ count $k_Y=1$.
The indices have lcm $n$, and
$\sqrt6\le1/n+\sum_{Y\in H}1/k_Y$. By Lemmas~\ref{blk:support}
and \ref{blk:height} the principal $3$-block has norm divisible by three,
contains the vacuum, and its supported nonunit indices are divisible
by three. Thus at least two of the five indices are multiples of three.

If $m=0$, the reciprocal sum is at most $3/2+2/3$. If $m=1$ and
at most one index is two, the four proper indices must be $(2,3,3,3)$,
of lcm six: the next upper bound $1/2+2/3+1/4$ falls short for $n\ge36$,
and at $n=30$ the last term is at most $1/5$. If two indices are two,
the others are multiples of three. An index three would leave a final
index $k$ with $n/k=\lcm(2,3,k)/k\le2$. Otherwise the reciprocal sum
is at most $2+2/6$.

If $m\ge4$ all supported degrees share a factor. For $m=3$, the two
proper degrees are coprime, hence at most $n/6$. The sum of entries of
a fixed degree is an integer by Galois permutation, so orthogonality
forces the three full-degree entries to sum to zero. A real row is
therefore impossible. A nonreal row is orthogonal to its conjugate,
so the six nonzero squares must give two of each cube root of unity.
The vacuum and the two singleton proper degrees already give three ones.

It remains that $m=2$. Repeated proper degrees have the form $a,a,b$
with $a,b$ coprime, contributing at most $2/b+1/a$ to the reciprocal
sum. If one is six the other is at least $35$; otherwise both are at
least ten. Either way $2+2/b+1/a+1/n<\sqrt6$. For three distinct
proper degrees their entries are rational signs, since each degree
occurs once. Together with the vacuum they exclude a nonreal row as above.
The two full signs cancel by orthogonality, since the proper degrees
sum to at most $3n/2$. Their indices $a<b<c$ thus satisfy
\[
 |n/a-n/b-n/c|=1,\qquad n=\lcm(a,b,c),\qquad
 (2n+1+n/a+n/b+n/c)^2\ge6n^2 .
\]
The equations and inequality give $2/a+2/n\ge\sqrt6-2$, hence $a\le5$.
Also $1/a\le2/b+1/n\le13/(6b)$ gives $b\le10$, and
$c(1/a-1/b)\le7/6$ gives $c\le7a(a+1)/6\le35$.
The remaining check is short enough to do by hand. For each $a,b$ let
$L=\lcm(a,b)$ and $t=(b-a)/\gcd(a,b)$. The signed equation becomes
$|tc-L|=\gcd(L,c)$, while $n/c\ge6$ gives $\gcd(L,c)\le L/6$.
Thus $5L/(6t)\le c\le7L/(6t)$, and at least two of $a,b,c$ are
multiples of three. The only pairs $(a,b)$ with a $c$ in this interval
and $b<c\le35$ satisfying that divisibility requirement are
$(2,3),(3,4),(3,6),(4,6),(5,6),(5,9)$. Dropping $tc=L$ leaves
$(3,6,7),(5,6,27),(5,6,33),(5,9,12)$. For the last three its left side minus
$6n^2$ is $-14900,-26136,-7776$, respectively. The first gives $n=42$
and the prime degree seven.
\end{proof}

\begin{lemma}\label{prime-square}
If $f_X=p^2$ and $d_X=n$, the nonunit degrees at entries of nonzero
absolute field trace divide $n/\gcd(n,p)$ and have gcd one.
\end{lemma}
\begin{proof}
For $u=\chi_X(Y)$ put $F_Y=\QQ(u)$ and $t_Y=\Tr_{F_Y/\QQ}(u)$.
Averaging the $[K_X:\QQ]$ orthogonal conjugate rows and grouping
columns by Galois orbit gives
\[
 f_X=\sum_{[Y]}[K_X:F_Y][K_Y:F_Y]t_Y^2,\qquad
 0=\sum_{[Y]}d_Y[K_Y:F_Y]t_Y.
\]
The vacuum contributes positively in the first sum and one in the
second, so $|t_Y|<p$ for $Y\ne0$ and the indicated degrees have gcd one.
Since $d_Y\mid pn$, failure of $d_Y\mid n$ makes $u$ divisible by $p$
by lcm integrality. If $v_p(d_Y)=v_p(n)>0$, a supported $Y$ is
principal and $u=(d_Y/n)\chi_Y(X)\equiv0$ above $p$. In both cases
$p\mid t_Y$, hence $t_Y=0$.
\end{proof}

\begin{lemma}\label{nine}
Nine is not a formal codegree.
\end{lemma}
\begin{proof}
Suppose $D=9n^2$. Let $I$ be the labels of degree $n$, $m=|I|$, and
$P$ the principal $3$-block. The $I$-rows are supported on $P$,
whose degrees divide $n$ by Lemma~\ref{blk:height}.

If $3\mid n$, fix $X\in I$. A supported $Y$ with $v_3(d_Y)=v_3(n)$
has entry $u=\chi_X(Y)\in\QQ(\zeta_9)$ reducing to zero above three.
Its $e=[\QQ(u):\QQ]$ conjugates occupy distinct columns, at cost
$C=\Tr(u\bar u)$. This is a positive multiple of three and
$C\ge e\,3^{2/e}$ by the norm. For $e\ge3$ this gives $C\ge9$,
exceeding the nonunit norm eight. Thus $u\in\QQ(\zeta_3)$ and has
trace zero by Lemma~\ref{prime-square}. Its contribution to the twisted
sum cancels that of the dual, whose twist agrees. The remaining nonunit
degrees divide $n/3$, so \eqref{st} after this cancellation gives
$3\le1/n+8/3$.

Now $3\nmid n$. The degrees in Lemma~\ref{prime-square} have gcd one and
each has at least two prime divisors. Thus $n$ has at least three primes
and $n\ge70$. Each $I$-row has support exactly $P$, so $|P|\le9$.
Suppose $I\cap P$ empty. At most seven nonunit degrees in $P$ equal
$n/2$ by gcd one, and every smaller divisor is at most $n/4$. Hence
$\sum_{Y\in P}d_Y^2\le1+7n^2/4+n^2/16<2n^2$. This mass is at least
$mn^2$: apply $s$ to $(d_Y\theta_Y)_{Y\in P}$ extended by zero, and
project to $I$, using the vacuum modular relation. Thus $m=1$, and its
unique row is integral. The nonunit norm eight gives either two double
entries, one double and four signs, or eight signs. Two doubles would
have coprime degrees, whereas two nonunits of $P$ meet by their
nonzero residues. For one double the twisted bound forces all five
degrees to equal $n/2$: a smaller degree loses at least $n/4$ from the
maximum $3n$. For eight signs at least four degrees equal $n/2$, since
$3n/2+5n/4<3n-1$. Their signs agree by
$(n/d_Y)u_Y=\chi_Y(X)\equiv n\pmod3$ and the remaining degrees cannot
cancel them to within one. This proves $I\cap P\ne\varnothing$.

A label in the intersection meets every $3$-free label by principality,
so $P$ is exactly the $3$-free set. All its twists are pure $3$-powers
by Lemma~\ref{blk:admissible} on that row. A $3$-active degree is
$3$-free, so three occurs in no mixed twist, contrary to the next lemma.
\end{proof}

\begin{lemma}\label{isolated}
If $D=(3n)^2$ and $3\nmid n$, then three occurs in some mixed twist.
\end{lemma}
\begin{proof}
Suppose otherwise. Put $W=W_3$, let $V$ be the complementary Weyl and
$\epsilon=\gamma^3$ with the lift of Section~\ref{blk:section}. On a
$3$-active twist fiber the complementary upper unipotents are scalar;
their conjugates have the same scalars since the complementary group
preserves the fiber. Thus $U(WUW^{-1})U=W^{-1}$ gives $V=\epsilon I$
there. At an active $Y$ its row meets no complementary active column,
since the complementary Weyls preserve its line. Lemma~\ref{blk:admissible}
then gives $f_Y=9$. On the span of all $3$-local translates of these fibers
$V=\epsilon I$ by commutation. On its orthogonal complement the
$3$-unipotent is trivial, as are its conjugates, which generate the local
group. Hence $W=I$ there and $s=VW=V+\epsilon(W-I)$.

The vacuum $v_0$ is fixed by the local tori, since the dimension character
is fixed and $D$ is a square. For a product $W_J$ of local Weyls,
Galois conjugation multiplies $W_J$ on the right by a local torus
\cite[Theorem II and Remark 4.5]{DLNCongruence}, so $w_J=W_Jv_0$ is
rational. Since $W_J^2$ is the local torus at $-1$, $w_J=W_J^{-1}v_0$.
For a CRT exponent $k$ equal to one in $J$ and zero elsewhere, the
unipotent identity identifies $sT^ks^{-1}v_0$ with $w_J$ up to diagonal
phases. The $i0$ and $0i$ entries agree by duality, so the indicator
formula at $f_ie_i$ in Lemma~\ref{blk:local} makes $f_i|(w_J)_i|^2$
integral. Thus $(3n/d_i)(w_J)_i$ is an integer.

In fact $3Wv_0$ is integral: at three use $3\nmid n$ and integrality
of $3nWv_0$; elsewhere subtract the untwisted term in $sT^ks^{-1}v_0$.
The difference at $i$ is $\sum_Y(\theta_Y^k-1)\chi_Y(i)/f_Y$, with
denominators nine. Similarly $Vv_0$ is integral at three: a contributing
$Y$ is active at some other prime, so its $3$-local group is trivial on
its line by the same fiber argument and $3\nmid f_Y$. For a $3$-active
$Y$, supplied by Cauchy, $(Wv_0)_Y=s_{Y0}/\epsilon=1/(3\epsilon)$.
Rationality gives $\epsilon=\pm1$.

Put $j=3(1-W_{00})$ and $a_i=3\epsilon(Wv_0)_i$ for $i\ne0$. The
splitting of $sv_0$ gives $a_i\equiv d_i/n\pmod3$ and, at the vacuum,
$nj\equiv-\epsilon\pmod3$. Thus $W_{00}\ne0$, so $|a_i|<3$ by
unitarity. The nonzero $a_i$ occur exactly at the $3$-free degrees,
and $0\le j\le6$. For $z=(W-I)v_0$ we have $\|z\|^2=2j/3$ and
$\langle z,sv_0\rangle=\epsilon j/3$, using $s=\epsilon W$ on the
moving subspace. Removing the vacuum term and subtracting gives
\[
 \sum_{i\ne0}(a_i^2-(d_i/n)a_i)=j(3-j)-\epsilon j/n.
\]
Every term is nonnegative, since $3$-free degrees divide $n$. The
vacuum lies in the principal block and $W$ preserves it. Thus the
$3$-free labels, which occur in $Wv_0$, are principal, and their
nonunit degrees pairwise share a prime. Their gcd is one by
Lemma~\ref{prime-square} at a $3$-active row. Removing occurrences
of $n$ preserves the gcd, as all the other degrees divide $n$.
There remain at least three distinct proper degrees, with distinct
indices $n/d_i$ prime to three. Their terms sum to at least
$3-1/2-1/4-1/5=41/20$. Yet $n\ge70$ since it has at least three
primes other than three, and the right side is at most $2+2/n<41/20$.
\end{proof}

\section{The largest degree classes}\label{sec:large}
\begin{lemma}\label{lem:eight-support}
If a largest degree $x$ has codegree eight, at least three odd primes divide $x$.
\end{lemma}
\begin{proof}
Here $D=8x^2$ and $M=2x$. Row $X$ of degree $x$ is supported in the
principal $2$-block, whose degrees divide $x$ by Lemma~\ref{blk:height}.
The supported nonunit degrees have gcd one, proving the claim for odd $x$:
each has at least two odd prime factors. Let $x$ be even. Since $M$ has
three prime factors, $x\ge15$. Use $u_Y,c_Y$ as in \eqref{st}. At every
hit, $\chi_Y(X)=(x/d_Y)u_Y$ reduces to $x=0$ above two.

If the row were rational, its entries would be integers and all signs
would occur at degrees dividing $a=x/2$. The nonunit norm seven gives
one double entry and three signs, or seven signs. A double at degree
$x$ could not be balanced, since $2x-1>3a$. At degree at most $a$ it
would give $\sqrt8\,x\le1+5a$. With seven signs the two signed degree
sums differ by one, and the smaller is at least $(\sqrt8\,x-2)/2>5a/2$.
Thus three signs of one kind all have degree $a$ and exactly two of the
other four have degree $a$: one or fewer gives at most $5a/2$, and
three or more leaves no possible balance within one. The two remaining
proper divisors of $a$ then sum to $a-1$. At least one is at most $a/3$,
so $a\le6$, impossible for two proper divisors at least six.

By Lemma~\ref{conduct} the row field lies in $\QQ(\zeta_8)$ and has degree
$h=2$ or $4$. The normalized traces $A_Y=\Tr(u_Y)/h$ are integers
since the normalized traces of powers of $\zeta_8$ are integers.
Averaging the $h$ orthogonal rows gives $\sum_{Y\ne0} A_Y^2=8/h-1$ and
$\sum_{Y\ne0}d_YA_Y=-1$. Hence $h=2$ and there are three nonunit signs.
In an imaginary quadratic field the conjugate rows are dual with equal
twists. Averaging their modular relations before taking absolute values
would give $\sqrt8\,x\le1+\sum_{Y\ne0,\,A_Y\ne0}d_Y\le1+5x/2$,
since their signed sum forces a proper divisor. The field is therefore
$\QQ(\sqrt2)$. Write $u_Y=A_Y+B_Y\sqrt2$ with integral $B_Y$. The norm
gives $\sum B_Y^2=2$, and reduction above two puts the three signs at
degrees dividing $x/2$.

Suppose the row misses degree $x$. At $d_Y\le x/4$,
$d_Y\sqrt{c_Y}\le xc_Y/4$. At $x/3$ the rational part vanishes, so a
hit has $c_Y=2$ and the same bound holds. At $x/2$ use
$2\sqrt{c_Y}\le c_Y+1$. For $k$ hits of degree $x/2$, \eqref{st} gives
$\sqrt8\le1/x+(7+k)/4$, hence $k\ge5$. Each nonunit divisor of $x$
shares a prime with $x/2$, so gcd one needs at least two further hits.
All seven would have $c_Y=1$, whereas two $B_Y$ are nonzero.

Thus $X$ meets a column of degree $x$. Its row too is supported in the
principal block, so $X$ is principal by symmetry and meets every odd
degree. The degree-$x$ entry has rational part zero, and its conjugate
occurs in a second column of degree $x$, using both nonzero $B_Y$.
Every odd label therefore hits a rational coordinate. There are four
such nonzero coordinates including the vacuum, and the number of odd
degrees counted with multiplicity is a positive multiple of four by
$D\equiv0\pmod4$. The three nonunit rational coordinates thus give odd
divisors of $x$ of signed sum $-1$, each with at least two prime
factors. Their gcd is one, requiring three odd primes.
\end{proof}


\begin{corollary}
At codegree eight, a largest degree $x$ has multiplicity at most two
when $x$ is even. For odd $x$ it is at most four if at most three odd
primes divide $x$, as holds for $r\le31$.
\end{corollary}
\begin{proof}
For even $x$, all maximal rows restrict to the same four signs on the odd labels: the vacuum sign is positive and the other three have signed degree sum $-1$. Since each degree exceeds one, the largest must occur alone on one side and the signs are determined. The squared projection of each normalized row on their common sign line is $4/8$, so there are at most two. For odd $x$, missing its own class would give $\sqrt8\,x\le1+7x/3$, since each proper divisor is at most $x/3$ and $\sqrt{c_Y}\le c_Y$ on a hit. Thus $X$ is principal at two by symmetry, and meets every label of degree $x$. When $r\le31$, four odd primes in $x$ would already require at least $2(1+2+4+6+10)>31$ labels by Proposition~\ref{blk:packets}. In the asserted case there are therefore only three odd primes to consider. Any two nonunit hit degrees therefore share a prime, since each contains at least two. The gcd-one condition requires three proper hits in addition to the vacuum, leaving at most four maximal labels.
\end{proof}

\subsection{Large classes in square dimension}
Write $s=S/\sqrt D$ and let $P_p$ denote the principal $p$-block. For rows $E$ supported on columns $A$, the vacuum modular relation and unitarity give
\begin{equation}\label{sq:mass}
 \sum_{X\in E}d_X^2=\|\operatorname{pr}_E sT(d_Y)_Y\|^2
       \le\sum_{Y\in A}d_Y^2.
\end{equation}
Two orthogonal unit columns of squared norm $m/f$ on a common subset have inner product there of modulus at most $1-m/f$, by Cauchy--Schwarz on its complement. Equality requires proportional restrictions of equal norm on that complement.

\begin{lemma}\label{sq:saturation}
A nonunit degree $n$ of prime-power codegree $f=p^a$ has multiplicity at most $f-2$.
\end{lemma}
\begin{proof}
Suppose its class $I$ has size $m=f-1$, the largest allowed by the vacuum. Its rows are supported on $P_p$. By symmetry $X\in I\setminus P_p$ would miss $I$. The complement has squared mass $n^2$, so equality in the vacuum modular relation gives $|s_{XY}|=d_Y/n$ there. At the vacuum this gives $n=\sqrt f$, contrary to the two prime divisors of $n$. Thus $I\subset P_p$.

Choose a nonunit $Y$ of degree $y$ prime to $p$. Each $\chi_X(Y)$ for $X\in I$ reduces to $y$ and is nonzero. Their squared moduli sum to a Galois-invariant integer, at least $m$ by the arithmetic--geometric mean on conjugates, and less than $f$ by column normalization and the vacuum. Thus both columns $Y,0$ have squared norm $m/f$ on $I$. Their inner product gives $|\sum_{X\in I}\chi_X(Y)|\le1$, strictly since their vacuum coordinates have different moduli. Yet this integer sum reduces to $my\ne0$ modulo $p$.
\end{proof}

\begin{proposition}\label{sq:fourteen}
If $D=M^2$ and $r\le31$, at most thirteen labels have degree $n=M/4$.
\end{proposition}
\begin{proof}
Lemma~\ref{sq:saturation} leaves only a class $I$ of fourteen. Put $P=P_2$. If $I\subset P$, the sums $w_Y=\sum_{X\in I}\chi_X(Y)$ are even integers. For nonunit $Y$ of odd degree the fourteen summands are nonzero. Their squared moduli sum to an even integer in $[14,16)$ by the same norm and residue arguments. The column bound gives $|w_Y|<2$, hence $w_Y=0$. This contradicts $\sum_{Y\ne0}d_Yw_Y/2=-7$ modulo two.

Thus $E=I\setminus P$ is nonempty, and $I\cap P$ is nonempty by~\eqref{sq:mass}. If $n$ were odd, a principal member of $I$ would meet all of $I$, forcing $I\subset P$. Hence $n$ is even. The rows of $E$ miss $I$ and see only dyadic twists. Some odd-primary twist is nontrivial outside $I$, whose complement has mass $2n^2$. By~\eqref{sq:mass}, $E$ consists of one Galois-fixed label $Z$, with integral row $u_Y=\chi_Z(Y)$.

At least six labels have a nontrivial odd twist component by Proposition~\ref{blk:packets} at two distinct odd primes. They are outside $I$, the vacuum and the support of $u$. If $u$ has at least nine nonunit hits, these, the vacuum and the thirteen labels of $I\cap P$ all belong to $P$ and are disjoint. The dyadic block of $Z$, intersected with the trivial odd-phase fiber, then has at most $31-6-9-1-13=2$ labels. Its span is dyadically invariant. A single label contradicts Lemma~\ref{blk:admissible}; two have equal degree by Lemma~\ref{blk:two}, whereas $Z$ is the only member of $I$ outside $P$.

It remains to find nine hits. At a nonunit hit $b_Y=M/d_Y$ satisfies $b_Yu_Y/4=\chi_Y(Z)\in2\mathbb Z$ by principality. As $1\le|u_Y|\le3$, we have $h_Y=b_Y/4=n/d_Y\ge2$ integral, and even if $u_Y$ is odd. The nonunit hits satisfy
\[
 \sum u_Y^2=15,\qquad \sum u_Y/h_Y=-1/n,\qquad
 \sum |u_Y|/h_Y\ge4-1/n.
\]
With at most eight hits the absolute weights are $(3,2,1,1)$, $(3,1^6)$ or $(2^3,1^3)$. The first has total weight seven, too small. In the others expand each weight into unit terms $1/h_Y$. Both signs have sum at least $2-1/n$, forcing one side to contain four terms and the other five.

If $n>42$, all four denominators on the first side are two. The other five terms must sum to $2\pm1/n$. Four or five twos overshoot, as each term exceeds $1/n$. With three twos the other two reciprocals sum to $1/2$ or differ from it by at least $1/42$ (the equalities are $(3,6),(4,4)$). With two twos the remaining three sum to one or at most $11/12$; with at most one two the total is at most $1/2+4/3$.

Finally an even $n\le42$ with two odd prime factors is $6q$, $q=5,7$. The only available proper degrees are $n/2,n/3,n/q$; odd weights require $n/2$. Thus $(3,1^6)$ violates gcd one. For $(2^3,1^3)$ two double weights must occur at $n/3,n/q$, giving total mass at most $9/2-2(1/2-1/3)-2(1/2-1/q)<4-1/n$.
\end{proof}

\begin{lemma}\label{sq:heights}
Suppose exactly three primes divide $M$. A denominator $p^a$ forces $p^2\mid M$, four denominators four force $8\mid M$, and a denominator five forces $125\mid M$. For each $\ell^e\Vert M$ there is a nonunit denominator divisible by $\lcm(\ell^e,p)$ in the first case, by $\lcm(\ell^e,4)$ in the second and by $\lcm(\ell^e,25)$ in the third.
\end{lemma}
\begin{proof}
A prime-power denominator row meets only $P_p$, hence only denominators divisible by $p$. The nonunit hit degrees have gcd one and each has at least two prime divisors. Their common multiple must therefore contain all three primes, which gives $p^2\mid M$.

For a four-row $X$ of degree $n=M/4$, suppose $b_Y\equiv2\bmod4$ at a hit. Then $\alpha=2s_{XY}$ is integral by lcm integrality. If $n$ is even, $(b_Y/2)\alpha=\chi_Y(X)$ vanishes modulo a prime above two by principality, so $|\operatorname{Norm}\alpha|\ge2$. Its $e$ distinct conjugates occur at distinct columns, contributing at least $e2^{2/e}/4\ge1$ before the vacuum. If $n$ is odd, each four-row is supported exactly on $P_2$ (by symmetry and principality of the columns there). The $Y$-column thus has at least four hits, each contributing at least $1/4$ after Galois averaging, again before the vacuum. Hence all hits have $4\mid b_Y$, giving $8\mid M$. For a five-row the nonzero trace degrees divide $(M/5)/\gcd(M/5,5)$ and have gcd one by Lemma~\ref{prime-square}; their common multiple again contains all three primes. This gives $125\mid M$ and $25\mid b_Y$ at each such hit. In each case a hit coprime to $\ell$ supplies the stated multiple.
\end{proof}

Consider a prefix of denominators with lcm $L$ and $h$ nonunit labels left. In square dimension every odd prime outside $L$ needs $p-1$ active suffix labels, so $p\le h+1$. If precisely three primes are possible, all must occur, and their sets of prime-free nonunit labels are disjoint. For each $p$, consider separately whether $v_p(M)=v_p(L)$ or is larger. In the latter case no prefix degree is $p$-free. In either case a $p$-free denominator must contain the full $p$-power in $M$, with the lower bounds of Lemma~\ref{sq:heights}. Reserve enough suffix positions for the two $p$-free nonunits of Lemma~\ref{lem:deletion}, and at least three at $2,5$ (reduce the dimension modulo $4,5$). Their counts at $2,3$ are respectively $3\bmod4$ and $2\bmod3$. Reserve a position also for each still missing witness in the lemma; several witnesses at one prime may use the same position, with the largest of their lower bounds. Round each lower bound $a_j$ up to an admissible multiple at least as large as the last prefix member. For $v$ reserved positions put $W=\sum_{j=1}^v\lceil L^2/a_j^2\rceil$ and $E=L^2(1-\sum b_i^{-2})$. Since the vacuum contributes at most one in units of $L^{-2}$, the next $b$ satisfies
\[
 b^2(E-W-1)\le (h-v)L^2 .
\]
Discard choices with $v>h$ and keep the weakest bound over feasible heights.

For the packet count at any $q$, set $F=B(q)$ and $j=(p-1)/2$ for each odd $p\mid M$, $p\nmid q$. A closed run of $m$ eligible equal degrees can contribute at most $j\lfloor m/j\rfloor$ labels towards $F(p-1)$; allow all $m$ in an open run. Writing $w$ for the sum of the positive deficits, we must have $w\le h$. Apply Proposition~\ref{blk:packets} to the deficient primes and bound the already present labels by the sum of their capacities. Adding a new odd prime $p\nmid qL$ to that union gives $w+F(p-1)\le h$, sharpening the bound for new primes.

\begin{proposition}\label{sq:mixed}
In square dimension the multiplicities $(i,j)$ of denominators four and five cannot be $(12,6),(10,9)$ or $(13,4)$. The pair $(13,3)$ excludes a denominator six.
\end{proposition}
\begin{proof}
For the first three pairs the remaining fraction of $D$ is $t/400$ with $t=4,6,11$, respectively. It is less than $1/25$, so the indicated multiplicities are already exact. Here $20\mid M$ and at least three primes divide $M$. For $t=4,6$ the only possibility with $M\le60$ is $60$, leaving nonunit squared mass $35$ or $53$, neither a sum of squares of degrees at least six. For $t=11$ we have $M\ge420$, and $M>1000$ if $25\mid M$: three-prime support forces $125\mid M$ by Lemma~\ref{sq:heights}, while four primes give $M\ge4\cdot3\cdot5\cdot7$, five times larger if $25\mid M$.

Let $I,J$ denote the two classes and $A,B$ their twist sums. The twists are dyadic and quinary, respectively, by Lemma~\ref{blk:admissible}. Independent conjugation of their fields and the Gauss sum bound give $|25\sigma A+16\tau B|\le t+400/M<13$. The variance of the conjugates of $A$ is integral: the dyadic cyclotomic power basis is orthonormal for normalized trace. The disk of radius $13/25$ makes that variance zero. Square Galois symmetry puts $B$ in $\mathbb Q(\sqrt5)$, and its conjugates differ by less than $26/16<\sqrt5$, so $B$ too is a rational integer. Now $A\equiv i\bmod2$, $B\equiv j\bmod5$, and $-j/4\le B\le j$ by the mean of a quinary root. The possible $(A,B)$ are $(-4,6)$, $(-6,9)$, and $(-3,4)$ or $(1,-1)$, respectively. Thus $25A+16B=\pm t$ except in the last case $(1,-1)$.

In an equality case write $25A+16B=\varepsilon t$. Averaging real parts of the Gauss sum gives
\[
 \sum_{Y\notin I\cup J}d_Y^2\bigl(1+\operatorname{mean}(\varepsilon\theta_Y)\bigr)\le M.
\]
For $p\mid M$, $p\ne2,5$, every $p$-active label is in this complement and its signed twist has mean at least $-1/2$: the root $-\varepsilon\theta_Y$ is nontrivial and has mean at most $1/2$. Thus $V_p\le2M$ and Lemma~\ref{blk:activity} gives smallest nonunit degree at least $M/6$. Every nonunit denominator is then at most six and divides $60$. The reciprocal-square identity multiplied by $60^2$ forces $M\mid60$, impossible.

In the exceptional case the twists on $J$ are the four primitive fifth roots: higher square orbits have at least ten members, and the only rational relation among the five fifth roots is their sum. Let $H$ be the labels outside $I\cup J$ of five-power twist, including the vacuum, with squared mass $\eta D$. The two classes contribute $9D/400$ to the Gauss sum. Their complement has mass $11D/400$, so
\[
 \sum_{Y\notin I\cup J}d_Y^2(1+\operatorname{mean}\theta_Y)\le D/200+M,
 \qquad \eta\le\tfrac43(1/200+1/M).
\]
Every $J$-row sees only five-power twists and is supported on $P_5$, hence misses $I$. One outside $P_5$ would also miss $J$ by symmetry, contrary to~\eqref{sq:mass} since $1/25>11/400$. All four are principal.

Put $y=M/5$, $\zeta=\zeta_5$. Their character fields lie in $\mathbb Q(\zeta_{25})$ by Lemma~\ref{conduct}, and hence in $\mathbb Q(\zeta)$ since their degrees are at most four. Their modular relations put the dyadic central charge $\xi$ in a five-power cyclotomic field, so $\xi=\pm1$. For $X\in J$ put
\[
 R_X=5\xi\theta_X^{-1}-\sum_{Z\in J}\theta_Z\chi_X(Z)
    =\frac{25}{M}\sum_{Z\in H}s_{XZ}\theta_Zd_Z\in\mathbb Z[\zeta].
\]
Unitarity at each conjugate gives $\sum_{X\in J}\operatorname{mean}|R_X|^2\le625\eta$. Galois covariance on both indices of the sum over $J$ gives $R_{\sigma X}=\sigma^2R_X$ for five-inertia.

If $5\mid y$, principality puts $R_X$ in $(1-\zeta)$. A nonzero element of this ideal has mean squared modulus at least $5/2$: write it as $\sum_{k=0}^4c_k\zeta^k$ with $\sum c_k=0$ by subtracting a common integer from the coefficients. The mean is $(5/4)\sum c_k^2$. Each inertia orbit in $J$ has at least two members by its primitive twists, whereas $625\eta\le(2500/3)(1/200+1/M)<5$. Hence all $R_X$ vanish. Then the projection of $sT(\mathbf d\,1_J)$ on $J$ already has its full squared norm $4y^2$ by the modular relation. Its vacuum coordinate $y^2B/M$ is nevertheless nonzero.

If $5\nmid y$, principality makes every five-free degree meet a $J$-row, hence have five-power twist. All five-active labels have five-free degree by Lemma~\ref{blk:admissible} and $v_5(M)=1$. On each nontrivial five-phase eigenspace every other local upper unipotent is scalar. Its local group preserves the space and is generated by conjugates of that unipotent; the local Weyl thus fixes each coordinate line. The same lemma makes each active codegree a five-power, hence its denominator is five. All twists on $H$ are therefore one and $\eta\le(1/200+1/M)/2$. But $R_X\equiv-4y\ne0\bmod(1-\zeta)$; each of the four terms has mean squared modulus at least one, contradicting $625\eta<4$ for $M\ge420$.

Finally $(13,3)$ together with a six leaves at most $143/3600<1/25$ for the set of denominators divisible by five other than five, including the vacuum. If $25\mid M$, the three $J$-rows have a rational member since their orbit lengths divide $[\mathbb Q(\zeta_{25}):\mathbb Q]=20$. Its nonzero entries have nonzero trace, so Lemma~\ref{prime-square} puts its support in that set, contrary to~\eqref{sq:mass}. If $25\nmid M$, the same mass bound again forces all $J$-rows to be principal. As just proved, this confines the five-active labels to $J$, whereas Proposition~\ref{blk:packets} requires four.
\end{proof}

\subsection{Ternary blocks in nonsquare dimension}
For $D=3M^2$, defect one at three means $v_3(f_X)=1$, or $3\nmid b_X$.
\begin{lemma}\label{ns:ternary}
In each local three-block containing such a label, the defect-one degrees are $(u,u,u)$ or $(u,u,2u)$; their denominators are $(z,z,z)$ or $(z,2z,2z)$.
\end{lemma}
\begin{proof}
Complete the fusion ring at three and extend its unramified base $\mathcal O$ to split residue fields. The generic factors $E_i$ of a block are totally ramified cyclotomic subfields, with ramification index $e_i$ and different exponent $\delta_i$. Coordinates in the order have the same residue. Its coefficient-of-unit trace $\sum_i\Tr_i/f_i$ makes it self-dual; the different bound of Lemma~\ref{conduct} gives $\delta_i\le e_iv_3(f_i)$.

A defect-one factor $i$ has local conductor at most three by Lemma~\ref{conduct}. Thus $e_i=1$ or $2$, and another factor $j$ must occur because $e_i/f_i$, the putative trace of the block identity, is nonintegral. On $E_i$ normalized trace $\lambda_i=\Tr_i/e_i$ lifts the residue. On $E_j$ take a linear residue lift $\lambda_j=\Tr_j(c\,\cdot)$, for instance the constant coefficient in a uniformizer basis. Since $\lambda_j(\mathfrak m_j)\subset3\mathcal O$ and $\lambda_j(1)=1$, the inverse different gives $v_j(c)=e_j-\delta_j-1$. Now $(\lambda_i-\lambda_j)/3$ is integral on the order. Its representing element belongs to the order by self-duality and has coordinates $f_i/(3e_i),-f_jc/3$, and zeros elsewhere. The first is a unit; equal residues exclude any third factor and force $v_3(f_j)=(\delta_j+1)/e_j$.

For $e_j=u3^s$ ($u=1,2$, $s\ge1$), the different formula in the proof of Lemma~\ref{conduct} gives $(\delta_j+1)/e_j=s+1-(1-3^{-s})/2$, strictly between integers. Thus $j$ too is tame of defect one. The two inertia orbits have lengths $a,b\in\{1,2\}$. Their phases have order dividing three by Lemma~\ref{blk:admissible} and are constant on each orbit by square Galois symmetry. Lemma~\ref{blk:moments} puts one at the trivial phase and one at order three; their degrees $u,v$ satisfy $bv^2=2au^2$. The only possibilities are $(a,b,v/u)=(1,2,1),(2,1,2)$.
\end{proof}
To test these triples on sorted denominators, at $z$ first fill the slots owed by $z/2$. Each remaining copy starts a doubled triple or belongs to an equal triple. Only the remainder modulo three need start doubled triples: three such triples can be replaced by one equal triple at $z$ and two at $2z$.

\section{The finite verification}\label{sec:verify}
Let $B$ be the least nonunit denominator, so $f=qB^2=D/n^2$ for
$n=\max_i d_i$. We have $f<r$ by the dimension equation. In fact
$f\le r-2$, since $f=r-1$ would give
$\sum_{i\ne0}(n^2-d_i^2)=1$, and each positive summand is at least
$2n-1$. Lemmas~\ref{small}, \ref{six} and~\ref{nine} give $f\ge8$,
excluding $r\le9$. Since $q$ is square-free and $f$ is composite,
every starting pair through rank $31$ occurs in this table:
\[
\begin{array}{c|l}
 B & q \\ \hline
 1 &10,14,15,21,22,26\\
 2 &2,3,5,6,7\\
 3 &2,3\\
 4,5 &1
\end{array}
\]

\subsection{Generating the lists}
The \supfile{square.cpp}{square} and \supfile{nonsquare.cpp}{nonsquare}
generators extend a nondecreasing denominator prefix by
\eqref{eq:interval}. With two nonunit positions left, they solve the
residual equation exactly. Branches exceeding native arithmetic pass
at their current prefix to the \supfile{square_exact.py}{square} or
\supfile{nonsquare_exact.py}{nonsquare} arbitrary-precision continuation.
In square dimension, if $k=M/L$ is even with three positions left,
the three remaining degree squares are all odd: together with the unit
they sum to $Ek^2\equiv0\pmod4$. Thus $Ek^2\equiv4\pmod8$,
so $v_2(k)=1$ and $E$ is odd.

For $q>1$ choose the full prime support $S$ of $M$ first. It contains
the primes of $qB$, has at least three members (four at codegree eight),
and satisfies the packet bound
\[
 B(q)\left(1+\sum_{p\in S,\,p\nmid2q}(p-1)\right)\le r.
\]
Hence every prime of $S$ is at most $r$. The
\supfile{primes.py}{prime list} is determined by the rank, and all
subsequent denominators are $S$-smooth. A further multiplicity bound applies when $v_p(f)=1$ at
an odd prime $p$. A maximal label then has trivial $p$-phase or order $p$.
The odd-valuation slice identity bounds their respective counts by $f/p$
and $(p-1)f/p$: the two masses are in ratio $1:(p-1)$ and sum to at most $D$.
The vacuum makes the first bound strict. In the principal $p$-block the
same identity supplies an active label, whose codegree is divisible by
$p^2$, making the second bound strict. The active count and its upper
bound are multiples of $(p-1)/2$. In total there are at most
$f-1-(p-1)/2$ maximal labels.

There is one further test of phase compatibility when $q=3$. The packet
bound gives $S=\{2,3,p\}$ with $p=5$ or $7$. A phase of order at least
$p^2$ would give at least $p(p-1)/2$ nontrivial phases, each with three
labels, besides the trivial fiber with at least three; already at $p=5$
this exceeds $31$. The nontrivial $p$-phases thus have order $p$ and
equal positive squared mass $W$. Take one fiber from each square class
of nonzero exponents modulo $p$. With degree counts $a_d,b_d$ and
multiplicities $m_d$, we need
\[
 \sum_d a_d d^2=\sum_d b_d d^2=W>0,\qquad
 (a_d+b_d)(p-1)/2\le m_d .
\]
Only nonunit $d$ with $p\mid D/d^2$ are eligible. Applying
\eqref{blk:fibergcd} at three to either $p$-phase fiber gives
$3g^2\mid W$, where $g$ is the gcd of its degrees.
The \supfile{screen.py}{terminal tests} compare vectors with the same $W$.

\subsection{Integer traces of character rows}\label{sq:orbits}
Let $X$ have degree $x>1$, codegree $f$ and Galois orbit of size
$h=[K_X:\mathbb Q]$, where $K_X$ is its character field.
The sum of its conjugate rows has integer coordinates
$z_Y=\Tr_{K_X/\mathbb Q}\chi_X(Y)$. Orthogonality gives
\begin{equation}\label{sq:orbit}
 z_0=h,\qquad \sum_{Y\ne0}d_Yz_Y=-h,\qquad
              \sum_{Y\ne0}z_Y^2=h(f-h).
\end{equation}
Thus $h\le\min(m_x,f-1)$. The $h$ coordinates on the orbit of $X$
coincide by the Galois action on arguments. If $h=1$, the vacuum
modular relation also gives $1+\sum_{Y\ne0}d_Y|z_Y|\ge\sqrt D$.

For $y=d_Y$ put $t=y/\gcd(x,y)$ and $F=\mathbb Q(\chi_X(Y))$.
By symmetry $F\subseteq K_X\cap K_Y$, and $\chi_X(Y)/t$ is an
algebraic integer. Since $[F:\mathbb Q]$ divides $h$ and
also divides $[K_Y:\mathbb Q]\le m_y$, the coordinate $z_Y$ is divisible by
$t h/\gcd(h,\lcm(1,\ldots,m_y))$. Coprime nonunit degrees have
zero crossing. If $f$ is a power of $p$, a hit requires $p^2\mid f_Y$
by Lemmas~\ref{blk:support} and~\ref{blk:height}; when $p\mid x$,
$(x/y)z_Y=\sum_{X'\in GX}\chi_Y(X')\in p\mathbb Z$ by principality.
At $f=p^2$, a nonzero coordinate requires $y\mid x/\gcd(p,x)$
by Lemma~\ref{prime-square}. The eligible nonunit degrees must have
gcd one: summing $\chi_X$ on each orbit of arguments gives the integer
$[K_Y:F]\Tr_{F/\mathbb Q}\chi_X(Y)$, and the degree-weighted sum is $-1$.
The \supfile{trace.py}{orbit test} solves \eqref{sq:orbit} for each $h$
at codegrees $f\le36$. A degree class must be a disjoint union of orbits
of feasible sizes.

\subsection{Tensor products}\label{sec:tensor}
Let $H$ be the additive monoid generated by the nonunit degrees.
The decomposition of $X\otimes X^*$ gives $a^2-1\in H$ for every
nonunit degree $a$. Fix a simple of degree $a$. Its fusion graph is
connected since it generates the category, allowing tensor words in
the object and its dual. Its unit vertex has only degree-$a$ neighbors.
Joining two nonunit degrees $u,v$ requires both $au-v$ and $av-u$
to be in $H$, with an optional unit subtracted from the first if $u=a$
and from the second if $v=a$, by reciprocity. All nonunit degrees must
be reachable from $a$ by such edges. In particular adjacent distinct
degrees in increasing order have ratio at most the least nonunit degree.

For a stronger test choose $X$ of least nonunit degree $x$. Let $C_{ab}$
sum the multiplicities of all degree-$b$ constituents of $X\otimes U$
over all degree-$a$ simples $U$. Reciprocity and braiding give
$N_{XU}^V=N_{XV^*}^{U^*}$, so $C_{ab}=C_{ba}$. Taking dimensions gives
\[
  \sum_{b>1}bC_{ab}=xa\,m_a-\mathbf1_{a=x},\qquad
  0\le C_{ab}\le m_am_b
       \lfloor\min(xa/b,xb/a,ab/x)\rfloor .
\]
The unit occurs just for $U=X^*$. The three bounds follow by rotating
the same multiplicity among the three products.
The \supfile{multiplication.py}{multiplication test} removes values
inconsistent with a row sum and shares the resulting bounds with the
transposed row. It also checks each individual product dimension, using
at degree $b$ at most the aggregate bound and at most $m_b$ times the
single-pair bound.

For example, suppose the degrees below $75$ are $1,6,10,12,12,15$.
The unique degree-$6$ product with itself gives
$35=2\cdot10+15$ after removing the unit. Thus $C_{10,6}=2$ and
$C_{12,6}=0$. The degree-$10$ product leaves $48=10i+12j+15k$
after removing its two sixes. Reduction modulo five and $j\le4$ give
$j=4$. Each degree-$12$ product has dimension $72<75$; together these
two products would give $144=4\cdot10+12i+15k$, impossible modulo
three. This excludes any completion of these small degrees.

\subsection{The dyadic trace test}\label{term:dyadic}
Suppose $D=16x^2$ and $m$ simples have the largest degree $x$.
When $m>8$, we may choose such an $X$ in the principal dyadic block
$P$. Otherwise all $m$ rows would miss the whole class by
Lemma~\ref{blk:support}; \eqref{sq:mass} would give $m\le16-m$.

For $y=d_Y$, put $g=\gcd(x,y)$, $t=y/g$, and
$\alpha_Y=\chi_X(Y)/t\in\mathbb Z[\zeta_{16}]$ by lcm integrality
and Lemma~\ref{conduct}. The power basis of this dyadic cyclotomic ring
is orthonormal for the normalized Hermitian trace. Thus the mean square
$C_Y$ of $\alpha_Y$ is the sum of the squares of its coefficients, and
its mean is its constant coefficient $k_Y$.

Only $y>1$ with $g>1$ and $4\mid f_Y$ can contribute off the vacuum.
In such a class of $s$ labels let $T_y=\sum k_Y$ and $B_y=\sum C_Y$.
Galois stability gives $\sum\alpha_Y=T_y$. If $x$ is even and $x/g$
odd, each $\alpha_Y$ reduces to $\delta=0$ above two: off $P$ it
vanishes, while on $P$ we have $(x/g)\alpha_Y=\chi_Y(X)\equiv x$.
If $X\in P$ and $t$ is odd, we may use $\delta=y\bmod2$ since
$t\alpha_Y\equiv y$. A prescribed residue gives
$T_y\equiv B_y\equiv s\delta\pmod2$ and $C_Y\ge F(k_Y)$, where
$F(k)=k^2+\mathbf1_{k\not\equiv\delta\ (2)}$: a parity mismatch
requires an odd nonconstant coefficient. Otherwise put $F(k)=k^2$.
Discrete convexity gives $B_y\ge(s-i)F(j)+iF(j+1)$ for
$|T_y|=sj+i$, $0\le i<s$.

The row equations and the averaged vacuum modular relation require
\[
 \sum_y ytT_y=-1,\qquad \sum_y t^2B_y=15,\qquad
 4x-1\le\sum_y yt\min\{\sqrt{sB_y}, B_y/\sqrt{\eta_y}\},
\]
where $t,s$ depend on $y$ and $\eta_y=2$ if $\delta=0$, otherwise $1$.
The two upper bounds follow from Cauchy--Schwarz and from
$C_Y=0$ or $C_Y\ge\eta_y$. The \supfile{screen.py}{dyadic test}
enumerates the fifteen available norm units and bounds each radical
term from above in units of $1/256$.

\clearpage
\subsection{Completion of the proof}\label{sec:finish}
The \supfile{verify.py}{verifier} generates the lists before applying
the terminal tests. The middle columns count complete degree lists at
that stage. The first row records the common result for each indicated
rank; rank nine also follows from the analytic bound above.
\[
\begin{array}{c|r@{\qquad}r|r}
 r & D\text{ square}&D\text{ nonsquare}&\text{survivors}\\ \hline
 9\text{--}24 &0&0&0\\
 25&1&0&0\\
 26&3&0&0\\
 27&10&0&0\\
 28&75&1&0\\
 29&829&1&0\\
 30&17\,165&19&0\\
 31&434\,509&428&0\\ \hline
 \text{total}&452\,592&449&0
\end{array}
\]
At each rank from $9$ to $24$ both generators exhaust their prefixes
without producing a complete list. Across the entire table, only three lists reach the
multiplication test, which excludes each of them. The degree-$6$
calculation in Section~\ref{sec:tensor} illustrates one of these exclusions.
\begin{proof}[Proof of the theorem]
A category of rank at most $31$ supplies one starting pair $(q,B)$.
Its ordered denominators follow a unique path in \eqref{eq:interval}.
Every discarded prefix violates a necessary condition, so its list
must reach the terminal tests and survive them. The last column
therefore proves the theorem.
\end{proof}
In the \supfile{README.md}{supplement directory}, run
\texttt{python3 verify.py --jobs 4} to regenerate the table. The
\supfile{certificate.json.gz}{certificate} records the source hashes,
partition of the search, and counts of profiles and exclusions for every
job. The command \texttt{python3 verify.py --check certificate.json.gz}
checks the partition and tallies against those records.
\begin{thebibliography}{9}
\bibitem{BNRWFinite}
P.~Bruillard, S.-H.~Ng, E.~C.~Rowell, and Z.~Wang,
\href{https://doi.org/10.1090/jams/842}{\emph{Rank-finiteness for modular categories}},
J. Amer. Math. Soc. \textbf{29} (2016), no.~3, 857--881.
\bibitem{BDP} S. Burciu, J. Dong, and S. Palcoux, \emph{Prime factors of fusion-category dimensions: weak group-theoreticity, rank and decomposition}, Hexagon 2611.00233 (2026). \url{https://hexagonmath.org/2611.00233}.

\bibitem{DLNCongruence}
C.~Dong, X.~Lin, and S.-H.~Ng,
\href{https://doi.org/10.2140/ant.2015.9.2121}{\emph{Congruence property in conformal field theory}},
Algebra \& Number Theory \textbf{9} (2015), no.~9, 2121--2166.
\bibitem{ENOannals}
P.~Etingof, D.~Nikshych, and V.~Ostrik,
\href{https://doi.org/10.4007/annals.2005.162.581}{\emph{On fusion categories}},
Ann. of Math. \textbf{162} (2005), no.~2, 581--642.
\bibitem{NSfs}
S.-H. Ng and P. Schauenburg, \href{https://doi.org/10.1016/j.aim.2006.07.017}{\emph{Frobenius--Schur indicators and exponents of spherical categories}}, Adv. Math. \textbf{211} (2007), 34--71.
\bibitem{NSind}
S.-H. Ng and P. Schauenburg, \href{https://doi.org/10.1007/s00220-010-1096-6}{\emph{Congruence subgroups and generalized Frobenius--Schur indicators}}, Comm. Math. Phys. 300 (2010), 1--46.
\end{thebibliography}
\end{document}
